\documentclass[11pt,letterpaper]{article}

\usepackage[utf8]{inputenc}
\usepackage[T1]{fontenc}
\usepackage{lmodern}
\usepackage{microtype}
\usepackage[margin=0.9in]{geometry}
\usepackage{amsmath, amssymb}
\usepackage{mathtools}
\usepackage{booktabs}
\usepackage{xcolor}
\usepackage[hidelinks]{hyperref}
\usepackage{enumitem}
\setlist{nosep, leftmargin=*}

\title{Boundary BC discretization --- implementation reference}
\author{\texttt{Research/Active/boundary\_conduction\_speedup/}}
\date{2026-04-29}

\begin{document}
\maketitle

\noindent\textbf{Scope.} Drop-in math reference for the three boundary modes
added to \texttt{Monodomain/Engine\_V5.4/cardiac\_sim/simulation/classical/discretization\_scheme/fdm.py}.
Keep open while implementing PLAN.md Phase~A.

\section{Setup}

5-point Laplacian on uniform $h$ grid. Indices: $i = 0, \dots, N_x-1$ along the
propagation axis $x$; $j = 0, \dots, N_y-1$ along $y$. Top wall at $j = 0$,
ghost row at $j = -1$.
\begin{equation}
\nabla^2 V \big|_{(i,j)} = \frac{V_{i-1,j} + V_{i+1,j} + V_{i,j-1} + V_{i,j+1} - 4 V_{i,j}}{h^2}
\end{equation}
Continuum BC: $\partial V / \partial y \big|_{y=0} = 0$.

\section{Three ghost-cell choices}

\begin{align}
\text{(A) face-centered, wall at } y &= -h/2 :   & \boxed{V_{i,-1} = V_{i,0}}     && \text{(V5.4 default since 2026-04-29)} \\
\text{(B) node-centered, wall at } y &= 0      :  & \boxed{V_{i,-1} = V_{i,1}}     && \text{(legacy V5.4 pre-flip)} \\
\text{(C) zero-pad,      no Neumann}             & : & \boxed{V_{i,-1} = 0}           && \text{(non-physical, diagnostic only)}
\end{align}

\section{Boundary stencil}

Decompose $\nabla^2 V = L_x + L_y$ with $L_x = V_{i-1,j}+V_{i+1,j}-2V_{i,j}$
and $L_y = V_{i,j-1}+V_{i,j+1}-2V_{i,j}$. At $j=0$, only $L_y$ depends on
the BC. Define $\varepsilon \coloneqq V_{i,1} - V_{i,0}$.
%
\begin{table}[h]\centering
\renewcommand{\arraystretch}{1.30}
\begin{tabular}{@{}lll@{}}
\toprule
\textbf{Mode} & $L_y \big|_{j=0}$ (substituting ghost) & In terms of $\varepsilon$ \\ \midrule
(A) face-centered  & $V_{i,0}+V_{i,1}-2V_{i,0}$  & $\varepsilon$ \\
(B) node-centered  & $V_{i,1}+V_{i,1}-2V_{i,0}$  & $2\varepsilon$ \\
(C) zero-pad       & $0    +V_{i,1}-2V_{i,0}$    & $\varepsilon - V_{i,0}$ \\
\bottomrule
\end{tabular}
\caption{Boundary $L_y$ under each BC. $L_x$ is identical across modes.}
\end{table}
%
\textbf{Implication.}
\begin{itemize}
\item Strict uniform-$y$ ($\varepsilon = 0$): A and B both give $L_y = 0$ (interior-equivalent); C gives $-V_{i,0}$ (always-on drain).
\item Non-uniform-$y$ ($\varepsilon \ne 0$): B amplifies any seed by exactly $2\times$ relative to A; C carries an additive $-V_{i,0}$ on top of A's $\varepsilon$.
\end{itemize}

\section{Hand-verification (corner)}

Test grid: $4 \times 4$, $V_{i,j} = 4i + j + 1$ (so $V_{2,3}=12$, $V_{3,2}=15$, $V_{3,3}=16$). Use $h=D=\chi=C_m=1$.
At the corner $(i,j)=(3,3)$ both east ($i+1 \ge N_x$) and north ($j+1 \ge N_y$) are off-grid; west $(2,3)=12$ and south $(3,2)=15$ are real.
%
\begin{table}[h]\centering
\renewcommand{\arraystretch}{1.30}
\begin{tabular}{@{}lll@{}}
\toprule
\textbf{Mode} & \textbf{Stencil sum} & $\nabla^2 V\big|_{(3,3)}$ \\ \midrule
(A) face-centered  & $V_{3,3} + V_{2,3} + V_{3,3} + V_{3,2} - 4V_{3,3} = 16+12+16+15-64$  & $-5$ \\
(B) node-centered  & $V_{2,3} + V_{2,3} + V_{3,2} + V_{3,2} - 4V_{3,3} = 12+12+15+15-64$  & $-10$ \\
(C) zero-pad       & $0       + V_{2,3} + 0       + V_{3,2} - 4V_{3,3} =     12+15-64$    & $-37$ \\
\bottomrule
\end{tabular}
\caption{Expected return values of \texttt{fdm.apply\_diffusion(V)[3,3]} for the three modes. Encoded as test assertions in PLAN.md Step~A.4.}
\end{table}

\section{Implementation patch (V5.4 fdm.py)}

Constructor change:
\begin{verbatim}
class FDMDiscretization(SpatialDiscretization):
    BOUNDARY_MODES = ('node_mirror_existing', 'face_mirror', 'zero_pad')

    def __init__(self, grid, D=1e-3, chi=1400.0, Cm=1.0, D_field=None,
                 boundary_mode='node_mirror_existing'):
        if boundary_mode not in self.BOUNDARY_MODES:
            raise ValueError(...)
        self.boundary_mode = boundary_mode
        ...
\end{verbatim}
%
At each off-grid neighbor in \texttt{\_build\_laplacian} (lines 264--320, four wall
branches: east \texttt{i+1 >= nx}, west \texttt{i-1 < 0}, north \texttt{j+1 >= ny},
south \texttt{j-1 < 0}; show the east branch, replicate symmetrically):
\begin{verbatim}
if self.boundary_mode == 'node_mirror_existing':
    # legacy: ghost = sub-edge.  Off-diagonal entry to (i-1,j) gets
    # +w from cardinal-west AND +w from this mirror, totalling 2w.
    A_rows.append(k); A_cols.append(idx(i-1, j)); A_vals.append(w_east)
    A_rows.append(k); A_cols.append(k);            A_vals.append(-w_east)
elif self.boundary_mode == 'face_mirror':
    # ghost = boundary cell itself.  East flux = D(V_C - V_C)/h^2 = 0.
    # Nothing added to off-diagonal; diagonal accumulator NOT incremented.
    pass
elif self.boundary_mode == 'zero_pad':
    # ghost = 0.  East flux = -D V_C / h^2.  Add to diagonal only.
    A_rows.append(k); A_cols.append(k); A_vals.append(-w_east)
\end{verbatim}
%
where \texttt{w\_east = D * cx} (use $D_{i,j}$ directly for zero-pad, since the ghost has no $D$; use the harmonic mean $D_\text{face}(D_{i,j}, D_{i\pm 1, j})$ for the mirror modes).

\section{Predicted behavior on a propagating wave}

\begin{table}[h]\centering
\renewcommand{\arraystretch}{1.25}
\begin{tabular}{@{}lll@{}}
\toprule
\textbf{Mode} & $\Delta(x_{\mathrm{mid}})$ \textbf{sign / mag.} & $h$\textbf{-scaling} \\ \midrule
(A) face-centered            & $\approx 0$           & flat at all $h$ \\
(B) node-centered (default)  & non-zero, sign TBD    & empirically determine $p$ in $|\Delta| \sim h^p$ \\
(C) zero-pad                 & non-zero, $> 0$       & empirically determine $p$ \\
bidomain bath-coupled        & $-7\%$ CV (Kleber)    & $\sim$ constant in $h$ (continuum) \\
\bottomrule
\end{tabular}
\caption{Predictions per H\textsubscript{BC}. Empirical $h$-scaling exponents
must be measured: the time-integral of $L_y$ along a propagating wavefront is
not closed-form in $h$ alone (depends on wave speed, upstroke shape, boundary-
layer length scale).}
\end{table}

\section{SPD compatibility}

\textbf{Correction (2026-04-29):} an earlier draft of this card claimed
``(B) writes a symmetric mirror entry'' --- this is wrong. Empirically
$\|L - L^\top\|_\infty = 1.0$ at the boundary for an $8\times8$ grid: at
$(N_x-1, j)$ the off-grid east neighbor mirrors to $(N_x-2, j)$ AND the
cardinal-west branch already writes to $(N_x-2, j)$, so after coalesce
$L[N_x-1, N_x-2] = +2w$. From $(N_x-2, j)$'s row, only the cardinal-east
branch writes back: $L[N_x-2, N_x-1] = +w$. The COO matrix is genuinely
asymmetric by $w$. The operator is self-adjoint in a weighted inner
product (boundary cells carry weight $1/2$), which is why PCG converges
even though the raw matrix is asymmetric --- but it is NOT a symmetric
matrix. Captured in \texttt{test\_boundary\_modes.py::test\_a3\_symmetry\_per\_mode}.

Symmetry status by mode:
\begin{itemize}
\item (A) \texttt{face\_mirror}: writes nothing to off-diagonal at the boundary; matrix is symmetric.
\item (B) \texttt{node\_mirror\_existing}: writes mirror off-diagonal that coalesces with the cardinal entry; matrix is asymmetric by $w$ at the boundary, but operator is self-adjoint in a weighted inner product.
\item (C) \texttt{zero\_pad}: only modifies diagonal at the boundary; matrix is symmetric.
\end{itemize}
All three are compatible with V5.4's PCG / Chebyshev solvers. The (B) mode
is also compatible with DCT / FFT solvers because PCG converges fine on
the asymmetric matrix --- but a strict ``symmetric Laplacian only''
implementation could legitimately reject it.

\textbf{V5.4 default flip (2026-04-29):} the default \texttt{boundary\_mode}
in \texttt{FDMDiscretization} was changed from \texttt{node\_mirror\_existing}
to \texttt{face\_mirror}. Reason: at the wall, $L_y\big|_{j=0}$ for node-mirror
is $2\varepsilon$ vs $\varepsilon$ for face-mirror, so node-mirror amplifies
any column-wise gradient by exactly $2\times$ at the boundary. This is the
root mechanism behind the storage-tank ``camel-toe'' / ``crescent''
artifacts that motivated this entire research question. face-mirror is the
genuine zero-flux Neumann choice (flux $\equiv 0$ identically, since
ghost = boundary cell makes the discrete face flux $D(V_{i,0}-V_{i,0})/h = 0$).
Anyone needing bit-exact reproduction of pre-flip simulations must pass
\texttt{boundary\_mode='node\_mirror\_existing'} explicitly.

The fourth conceivable mode --- node-mirror-with-fold, which routes outgoing flux into the sub-edge asymmetrically --- breaks symmetry harder. \emph{Deferred to Phase~F}; would require a non-symmetric solver (GMRES / BiCGStab) which V5.4 does not currently provide.

\end{document}
